\section{Prompt Processing, Science Platform and the APDB}\label{sec:upservices}

This replaces the Alert Production rows of \tabref{tab:drpAndAlertSizing} and
the US Data Access Centre and staff rows of \tabref{tab:lspSizing} and
\tabref{tab:computeSizingOps}. The Chilean Data Access Centre rows are not
updated. These resources are provided by the USDF alone: partner facilities
neither run Prompt Processing nor host the Data Access Centre. All of them
run on Kubernetes, not on the batch fleet, so none of them appears in the DRP
compute of \secref{sec:upcompute}.

\begin{itemize}
\item \emph{Prompt Processing}, which includes Alert Production and alert
      distribution, runs on Kubernetes. Its core count is an input held flat
      across the survey; its disk is in \secref{sec:upstorage}.
\item \emph{Data Access Centre} compute is 10\% of DRP peak demand, and
      \emph{staff} compute is 10\% of that. These are the same proportions as
      the user compute described in the construction budget, now applied to
      the measured DRP basis. As before, the Rubin Science Platform's compute
      is part of this allowance, not additional to it.
\item \emph{Interactive} nodes are counted as deployed.
\item The \emph{APDB} grows with cumulative visits, scaled from a reference
      size at a reference visit count. Its on-floor volume is inside the other
      and miscellaneous storage allowance (\secref{sec:upstorage}), not added
      to it. The APDB is also part of the fast tier. Its off-site backup
      accumulates to an end-of-survey figure.
\end{itemize}

The service line in \tabref{tab:smServices} compares Prompt Processing, Data
Access Centre and staff cores with the deployed Kubernetes fleet. It is a
headroom check, not a purchase trigger; Kubernetes growth is costed
separately (\secref{sec:upbudget}). One inconsistency remains open: planning
documents give 48 allocatable cores per Kubernetes node, but the deployed
nodes have 64 hardware threads (DM-54218).

% GENERATED by generate_model.py from lsst/rubin-sizing-model @ v2026.09.2. Do not edit by hand.
% Regenerate: git clone --branch v2026.09.2 https://github.com/lsst/rubin-sizing-model ; then, inside the clone, uv run generate_model.py --emit-tex <this directory>

\small \begin{longtable}{|p{0.55\textwidth}|r|l|}
\caption{Alert Production, user and staff compute, and APDB inputs (USDF only). \label{tab:smServiceInputs}}\\
\hline
\textbf{Metric} & \textbf{Value} & \textbf{Unit} \\
\hline\endfirsthead
\hline
\textbf{Metric} & \textbf{Value} & \textbf{Unit} \\
\hline\endhead
Alert Production cores & 1,188 & cores \\ \hline
Alerts per visit & 10,000 & alerts \\ \hline
Cores per K8s node & 48 & cores \\ \hline
K8s nodes deployed & 75 & nodes \\ \hline
DAC / LSP cores as fraction of DRP & 10\% & fraction \\ \hline
Staff cores as fraction of DAC & 10\% & fraction \\ \hline
Interactive nodes & 11 & nodes \\ \hline
Cores per interactive node & 64 & cores \\ \hline
APDB reference size & 4.5 & TB \\ \hline
APDB reference visits & 57,000 & visits \\ \hline
Cassandra nodes deployed & 12 & nodes \\ \hline
APDB off-site backup at end of survey & 100 & TB \\ \hline
\end{longtable} \normalsize


\begin{landscape}
% GENERATED by generate_model.py from lsst/rubin-sizing-model @ v2026.10.0. Do not edit by hand.
% Regenerate: git clone --branch v2026.10.0 https://github.com/lsst/rubin-sizing-model ; then, inside the clone, uv run generate_model.py --emit-tex <this directory>

\tiny \begin{longtable}{|p{0.22\textwidth}|l|r|r|r|r|r|r|r|r|r|r|}
\caption{Prompt Processing, user, staff and APDB resources at the USDF, by operations year. All run on Kubernetes. \label{tab:smServices}}\\
\hline
\textbf{Metric} & \textbf{Unit} & \textbf{Year 1 (DP2, actual)} & \textbf{Year 2} & \textbf{Year 3} & \textbf{Year 4} & \textbf{Year 5} & \textbf{Year 6} & \textbf{Year 7} & \textbf{Year 8} & \textbf{Year 9} & \textbf{Year 10} \\
 &  & FY2026 & FY2027 & FY2028 & FY2029 & FY2030 & FY2031 & FY2032 & FY2033 & FY2034 & FY2035 \\
\hline\endfirsthead
\hline
\textbf{Metric} & \textbf{Unit} & \textbf{Year 1 (DP2, actual)} & \textbf{Year 2} & \textbf{Year 3} & \textbf{Year 4} & \textbf{Year 5} & \textbf{Year 6} & \textbf{Year 7} & \textbf{Year 8} & \textbf{Year 9} & \textbf{Year 10} \\
 &  & FY2026 & FY2027 & FY2028 & FY2029 & FY2030 & FY2031 & FY2032 & FY2033 & FY2034 & FY2035 \\
\hline\endhead
Prompt Processing cores & cores & 1,188 & 1,188 & 1,188 & 1,188 & 1,188 & 1,188 & 1,188 & 1,188 & 1,188 & 1,188 \\ \hline
Prompt Processing K8s nodes & nodes & 25 & 25 & 25 & 25 & 25 & 25 & 25 & 25 & 25 & 25 \\ \hline
DAC / LSP cores (share of DRP peak) & cores & 79 & 1,008 & 2,016 & 3,024 & 4,032 & 5,040 & 6,047 & 7,055 & 8,063 & 9,071 \\ \hline
Staff cores & cores & 8 & 101 & 202 & 302 & 403 & 504 & 605 & 706 & 806 & 907 \\ \hline
Interactive nodes & nodes & 11 & 11 & 11 & 11 & 11 & 11 & 11 & 11 & 11 & 11 \\ \hline
Service cores vs deployed K8s fleet & \% & 35\% & 64\% & 95\% & 125\% & 156\% & 187\% & 218\% & 249\% & 279\% & 310\% \\ \hline
APDB on floor & TB & 16.6 & 33.2 & 49.7 & 66.3 & 82.9 & 99.5 & 116.1 & 132.6 & 149.2 & 165.8 \\ \hline
APDB off-site backup, cumulative & TB & 10 & 20 & 30 & 40 & 50 & 60 & 70 & 80 & 90 & 100 \\ \hline
Cassandra nodes deployed & nodes & 12 & 12 & 12 & 12 & 12 & 12 & 12 & 12 & 12 & 12 \\ \hline
\end{longtable} \normalsize

\end{landscape}
